\par\Needspace{4\baselineskip}
\begin{center}
\hypertarget{native-061-article044}{}\pdfbookmark[2]{44.}{article-044}44.
\end{center}
Die \emph{hyperbolische} Bewegung können wir kürzer behandeln, da das Verfahren dem bisher für die elliptische Bewegung erläuterten ganz analog ist. Wir schreiben die Gleichung zwischen der Zeit $t$ und der Hilfsgröße $u$ in folgender Form.
\[
\begin{aligned}
&(e-1)\left\{\tfrac{1}{20}\left(u-\frac{1}{u}\right)+\tfrac{9}{10}\log u\right\}\\
&\quad+\left(\tfrac{1}{10}+\tfrac{9}{10}e\right)\left\{\tfrac{1}{2}\left(u-\frac{1}{u}\right)-\log u\right\}=kt\left(\frac{e-1}{q}\right)^{\frac{3}{2}},
\end{aligned}
\]
\noindent Hier werden hyperbolische Logarithmen verwendet. Dabei ist $\tfrac{1}{20}\left(u-\dfrac{1}{u}\right)+\tfrac{9}{10}\log u$ eine Größe erster Ordnung und $\tfrac{1}{2}\left(u-\dfrac{1}{u}\right)-\log u$ eine Größe dritter Ordnung, wenn man $\log u$ als kleine Größe erster Ordnung auffasst. Setzen wir also
\[
\begin{aligned}
\frac{6\left\{\tfrac{1}{2}\left(u-\frac{1}{u}\right)-\log u\right\}}{\tfrac{1}{20}\left(u-\frac{1}{u}\right)+\tfrac{9}{10}\log u}&=4A,\\[.7em]
\frac{\tfrac{1}{20}\left(u-\frac{1}{u}\right)+\tfrac{9}{10}\log u}{2\BareRoot{A}}&=B,
\end{aligned}
\]
\noindent\SourceLeaf{062}so ist $A$ eine Größe zweiter Ordnung, während $B$ von der Einheit um eine Größe vierter Ordnung abweicht. Unsere Gleichung nimmt dann die folgende Form an.
\[
\text{2)}\qquad B\left\{2(e-1)A^{\frac{1}{2}}+\tfrac{2}{15}(1+9e)A^{\frac{3}{2}}\right\}=kt\left(\frac{e-1}{q}\right)^{\frac{3}{2}},
\]
\noindent Sie ist Gleichung 1) in Artikel 37 ganz analog. Setzen wir $\left(\dfrac{u-1}{u+1}\right)^2=T$, so ist $T$ von zweiter Ordnung, und mit der Methode der unendlichen Reihen ergibt sich
\[
\begin{aligned}
\frac{A}{T}&=1+\tfrac{4}{5}A+\tfrac{8}{175}A^2-\tfrac{8}{525}A^3\\
&\quad+\tfrac{1896}{336875}A^4-\tfrac{28744}{13138125}A^5+\text{usw.}
\end{aligned}
\]
\noindent Setzen wir daher $\dfrac{A}{T}=1+\tfrac{4}{5}A+C$, so ist $C$ eine Größe vierter Ordnung, und es gilt $A=\dfrac{(1+C)T}{1-\tfrac{4}{5}T}\mathord{\raisebox{.35ex}{.}}$ Für den Radiusvektor folgt schließlich leicht aus Gleichung VII in Art. 21
\[
r=\frac{q}{(1-T)\cos\tfrac{1}{2}v^2}=\frac{(1+\tfrac{4}{5}A+C)q}{(1-\tfrac{1}{5}A+C)\cos\tfrac{1}{2}v^2}\mathord{\raisebox{.35ex}{.}}
\]
